% ══════════════════════════════════════════════════════════════════════
% SECTION V
% ══════════════════════════════════════════════════════════════════════
\section{Operationalizing Agency: Economics, Safety, and Latency}
\label{sec:operational}

\subsection{The Unreliability Tax and the Cost of Autonomy}

Agentic AI introduces a fundamentally different cost model from traditional software infrastructure. Infrastructure costs are relatively fixed; intelligence costs are variable and correlated with agent complexity. When agents fail and retry, costs compound quickly.

Research into the \keyword{Unreliability Tax} identifies the additional compute, latency, and engineering required to compensate for agent failures~\citep{measurement_imbalance}. For instance, a ``Reflexion'' loop that runs for 10 cycles can consume \textbf{50 times the tokens} of a single linear pass. This is not a marginal increase---it represents a qualitative change in the economics of AI deployment.

The Unreliability Tax manifests in several ways:

\begin{itemize}
    \item \textbf{Retry costs}: When agents fail, they must retry, consuming additional tokens and API calls. In complex workflows, a single failure can trigger a cascade of retries across multiple steps.
    \item \textbf{Guardrail costs}: Implementing safety checks, validation steps, and human-in-the-loop checkpoints adds latency and token consumption to every workflow.
    \item \textbf{Observability costs}: Capturing and storing execution traces for evaluation and debugging requires infrastructure and storage.
    \item \textbf{Engineering costs}: Building and maintaining the evaluation infrastructure, red teaming pipelines, and monitoring systems represents a significant ongoing investment.
\end{itemize}

These costs are often invisible in benchmark evaluations, which typically measure accuracy in isolation. In production, they can be the difference between a viable product and an economic failure.

\subsection{CNA (Cost-Normalized Accuracy) as a New Industry Standard}

To evaluate economic efficiency, the metric of \keyword{Cost-Normalized Accuracy (CNA)} is used~\citep{beyond_accuracy}:
\begin{equation}
    \text{CNA} = \frac{\text{Accuracy}}{\text{Cost}} \times 100
\end{equation}
where Cost is measured in USD per task. CNA makes the accuracy-cost tradeoff explicit, enabling organizations to compare agents not just on how well they perform, but on how efficiently they achieve that performance.

Evaluation of leading agent architectures on enterprise tasks demonstrates a striking result, shown in Table~\ref{tab:cna_results} and visualized in Figure~\ref{fig:cna_comparison}.

\begin{table}[htbp]
\centering
\caption{Architecture comparison: efficacy, cost, and CNA.}
\label{tab:cna_results}
\resizebox{0.96\columnwidth}{!}{%
\begin{tabular}{@{}lccc@{}}
\toprule
\textbf{Architecture} & \textbf{Efficacy (\%)} & \textbf{Cost (USD/Task)} & \textbf{CNA} \\
\midrule
ReAct-GPT-o3 & 68.7 & 0.85 & 80.8 \\
Reflexion & 74.1 & 4.35 & 17.0 \\
Domain-Tuned & 81.5 & 0.31 & 260.4 \\
Plan-Execute & 71.9 & 1.05 & 68.5 \\
\bottomrule
\end{tabular}
}
\end{table}

\begin{figure*}[htbp]
\centering
% Cost-Normalized Accuracy Comparison Across Architectures
\begin{tikzpicture}
\begin{axis}[
    ybar,
    width=14cm, height=7cm,
    bar width=12pt,
    ylabel={Score (\%)},
    ylabel style={font=\small},
    symbolic x coords={ReAct, Reflexion, Domain-Tuned, Plan-Execute},
    xtick=data,
    xticklabel style={font=\small},
    yticklabel style={font=\small},
    ymin=0, ymax=100,
    legend style={at={(0.5,-0.18)}, anchor=north, legend columns=3, font=\small},
    enlarge x limits=0.2,
    nodes near coords,
    nodes near coords style={font=\tiny},
    every axis plot/.append style={fill opacity=0.85},
    grid=major,
    grid style={gray!20},
    title={\textbf{Architecture Comparison: Efficacy, Cost, and CNA}},
    title style={font=\normalsize, text=alchedark}
]

% Efficacy (task success rate)
\addplot[fill=pillar1!70, draw=pillar1!90] coordinates {
    (ReAct, 62) (Reflexion, 71) (Domain-Tuned, 78) (Plan-Execute, 74)
};

% Cost efficiency (inverse of relative cost, normalized)
\addplot[fill=pillar3!70, draw=pillar3!90] coordinates {
    (ReAct, 45) (Reflexion, 38) (Domain-Tuned, 65) (Plan-Execute, 52)
};

% CNA (cost-normalized accuracy)
\addplot[fill=pillar2!70, draw=pillar2!90] coordinates {
    (ReAct, 28) (Reflexion, 27) (Domain-Tuned, 51) (Plan-Execute, 38)
};

\legend{Efficacy, Cost Efficiency, CNA}
\end{axis}
\end{tikzpicture}
\caption{Architecture comparison across efficacy, cost efficiency, and Cost-Normalized Accuracy (CNA). The Domain-Tuned architecture achieves the highest CNA by combining strong efficacy with low cost, while Reflexion's high efficacy is undermined by excessive cost.}
\label{fig:cna_comparison}
\end{figure*}

The Reflexion architecture---often cited for its high accuracy---is \textbf{4.7$\times$ more expensive} than the domain-tuned alternative while delivering lower accuracy. The Domain-Tuned agent achieves the highest accuracy at the lowest cost, producing a CNA more than 15$\times$ higher than Reflexion.

This finding highlights what researchers call a \keyword{distorted research landscape} where expensive and fragile solutions appear superior in standard accuracy-only benchmarks~\citep{measurement_imbalance}. Simple baseline strategies can sometimes outperform complex agents at 50$\times$ lower cost.

\textbf{CNA should be a standard reporting requirement alongside any accuracy metric.} Any benchmark result reported without cost normalization is incomplete. The most accurate agent in an evaluation suite may be economically inviable at production scale.

\subsection{Security Challenges: Offensive Agentic Risk and API/MCP Gateways}

As agentic systems gain the ability to take autonomous actions, the security implications become critical~\citep{securing_ai_agents}. The attack surface of agentic environments maps across four layers, as illustrated in Figure~\ref{fig:security_layers}.

\begin{figure*}[htbp]
\centering
% Security Attack Surface: Four Layers of Agentic AI Vulnerability
\resizebox{0.96\textwidth}{!}{%
\begin{tikzpicture}[
    layer/.style={rectangle, rounded corners=6pt, draw=#1!70!black, fill=#1!12,
        minimum width=12cm, minimum height=1.3cm, align=center, font=\small,
        inner sep=6pt},
    threat/.style={rectangle, rounded corners=3pt, draw=riskred!60, fill=riskred!8,
        minimum width=2.8cm, minimum height=0.6cm, align=center, font=\scriptsize,
        text=riskred!80},
    arr/.style={-{Stealth[length=4pt]}, thick, alchegray!60}
]

% Layers (bottom to top)
\node[layer=pillar4] (l1) at (0, 0) {\textbf{Runtime Layer} \quad Container escapes, resource exhaustion, sandbox breaches};
\node[layer=pillar3] (l2) at (0, 1.8) {\textbf{SaaS / Application Layer} \quad Data exfiltration, privilege escalation, API abuse};
\node[layer=pillar2] (l3) at (0, 3.6) {\textbf{API / MCP Layer} \quad Prompt injection, tool poisoning, context manipulation};
\node[layer=pillar1] (l4) at (0, 5.4) {\textbf{Endpoint / User Layer} \quad Social engineering, credential theft, session hijacking};

% Arrows between layers
\draw[arr] (l1) -- (l2);
\draw[arr] (l2) -- (l3);
\draw[arr] (l3) -- (l4);

% Threat examples on the right
\node[threat] (t1) at (8.5, 0) {Container escape};
\node[threat] (t2) at (8.5, 1.8) {Data exfiltration};
\node[threat] (t3) at (8.5, 3.6) {Prompt injection};
\node[threat] (t4) at (8.5, 5.4) {Social engineering};

\draw[riskred!40, dashed, -{Stealth[length=3pt]}] (l1.east) -- (t1.west);
\draw[riskred!40, dashed, -{Stealth[length=3pt]}] (l2.east) -- (t2.west);
\draw[riskred!40, dashed, -{Stealth[length=3pt]}] (l3.east) -- (t3.west);
\draw[riskred!40, dashed, -{Stealth[length=3pt]}] (l4.east) -- (t4.west);

% Title
\node[font=\Large\bfseries, text=alchedark] at (0, 7.0) {Agentic AI Security Attack Surface};

% Left brace with label
\draw[decorate, decoration={brace, amplitude=8pt, mirror}, thick, alcheblue!60]
    (-6.9, -0.65) -- (-6.9, 6.05)
    node[midway, left=10pt, font=\small\bfseries, text=alcheblue, align=center] {Increasing\\attack\\surface};

\end{tikzpicture}
}
\caption{The four-layer security attack surface for agentic AI systems, from runtime infrastructure to endpoint users, with representative threat vectors at each layer.}
\label{fig:security_layers}
\end{figure*}

\begin{enumerate}
    \item \textbf{The Endpoint Layer} (e.g., coding agents): Compromise of the agent's execution environment, including local file systems, development environments, and developer machines.
    \item \textbf{The API/MCP Gateway Layer} (tool exchanges): Manipulation of the tool interfaces that agents use to interact with external systems. A compromised API can feed the agent malicious data, causing it to take harmful actions.
    \item \textbf{The SaaS Platform Layer}: Exploitation of the SaaS applications that agents interact with---CRM systems, communication platforms, cloud services---through the agent's authorized access.
    \item \textbf{The Runtime Layer}: Attacks on the agent's execution runtime, including prompt injection, context manipulation, and goal hijacking.
\end{enumerate}

CISOs are encouraged to implement a three-stage security framework:

\begin{itemize}
    \item \textbf{Visibility}: Knowing what agents you have, what tools they can access, and what actions they are taking. You cannot secure what you cannot see.
    \item \textbf{Configuration}: Reducing the blast radius by implementing least-privilege access, tool restrictions, and action boundaries. Agents should have the minimum permissions necessary for their tasks.
    \item \textbf{Runtime Protection}: Monitoring agent behavior in real-time for anomalous actions, policy violations, and indicators of compromise.
\end{itemize}

The security challenge is compounded by the fact that agentic systems can be manipulated through their inputs in ways that traditional software cannot. Prompt injection---embedding malicious instructions in data that the agent processes---can cause agents to take actions that violate their intended behavior. Tool poisoning---manipulating the outputs of tools that the agent relies on---can cause agents to make incorrect decisions based on falsified information.

These security considerations must be integrated into evaluation from the beginning, not bolted on after deployment. An evaluation framework that does not test for security vulnerabilities is providing a false sense of security.